% Springer LNCS manuscript skeleton for the Bangla LLM fragmentation/calibration study.
% All TODOs must be resolved from the final frozen experiment. Do not invent results.
% Compile with the official Springer LNCS llncs.cls and splncs04.bst.

\documentclass[runningheads]{llncs}

\usepackage[T1]{fontenc}
\usepackage{amsmath,amssymb}
\usepackage{graphicx}
\usepackage{booktabs}
\usepackage{multirow}
\usepackage{url}
\usepackage{hyperref}
\usepackage{xurl}
\usepackage{microtype}
\usepackage[utf8]{inputenc}


\hypersetup{colorlinks=true,linkcolor=black,citecolor=black,urlcolor=blue}
\urlstyle{rm}

\begin{document}

\title{Does Tokenization Fragmentation Predict Miscalibration in Bangla LLMs? A Cross-Lingual Study}
\titlerunning{Tokenization Fragmentation and LLM Calibration}

% Current working team: Subho, Naeem, and Abdullah.
% Confirm exact full names, author order, affiliations, emails, and ORCIDs with all three authors.
\author{Md. Razaul Haque Subho \and Naeem [Full Name to Confirm] \and Abdullah al Azmi}
\authorrunning{Subho et al.}
\institute{Affiliations to be confirmed by the authors.}
\institute{Affiliations to be confirmed by the authors.}
\institute{Department of Mechanical and Production Engineering \\ Islamic University of Technology (IUT), Gazipur, Bangladesh}

\maketitle

\begin{abstract}
	% Write this LAST, after results are frozen (approximately 150--250 words).
	% Include: problem, gap, research questions, dataset and model panel, measurement
	% protocol, principal quantitative findings with uncertainty, and contribution.
	% Do not include unverified or preliminary branch results.
\end{abstract}

\keywords{Large language models \and Bangla \and Tokenization \and Fragmentation \and
	Confidence calibration \and Cross-lingual evaluation}

\section{Introduction}
\subsection{Background and Motivation}
% Explain the use of LLMs for Bangla and why language representation matters.
% Introduce subword tokenization and confidence calibration in accessible terms.

\subsection{Problem Statement}
% State the problem precisely: whether item-level tokenization fragmentation is
% associated with model confidence/calibration behaviour on Bangla tasks.
% Distinguish accuracy (being correct) from calibration (confidence matching accuracy).

\subsection{Research Gap}
% Support the gap with the reviewed literature. Do not claim that no prior work exists
% unless the literature review justifies that statement.
% State what prior tokenizer studies and Bangla LLM evaluations do not jointly test.

\subsection{Research Questions and Contributions}
\begin{itemize}
	\item \textbf{RQ1:} How does tokenization fragmentation differ between Bangla and English under the selected model tokenizers?
	\item \textbf{RQ2:} Is item-level tokenization fragmentation associated with confidence and calibration behaviour on Bangla evaluation items?
	\item \textbf{RQ3:} Does this association vary across models, subjects, or language conditions?
\end{itemize}
% List 2--4 concrete contributions only after the protocol and evidence are established.

\subsection{Paper Organization}
% One short paragraph explaining the remaining sections.

\section{Related Work}
\subsection{Subword Tokenization and Multilingual Tokenizer Behaviour}
% Cover BPE, WordPiece, SentencePiece/unigram or relevant families as applicable.
% Explain vocabulary sharing, language coverage, and fragmentation.
% Cite primary papers and report their limitations accurately.

\subsection{Measuring Tokenization Fragmentation}
% Review TPW and other measures used in the literature.
% Explain what each measure captures and what it misses.
% STRR must not be included as a settled metric until its exact source definition
% and implementation have been checked.

\subsection{Bangla LLMs and Bangla Benchmarking}
% Review the selected Bangla-capable models and BnMMLU.
% Distinguish original Bangla benchmark items from translated English counterparts.

\subsection{Confidence Estimation and Calibration in LLMs}
% Review confidence extraction from multiple-choice logits/probabilities,
% ECE, Brier score, reliability diagrams, and relevant calibration work.

\subsection{Summary of the Gap}
% End the section by stating the specific gap addressed by this paper.

\section{Study Design and Research Questions}
\subsection{Study Overview}
% Describe the unit of comparison: model x item x language condition.
% Clarify which comparisons are paired (same source question translated)
% and which are unpaired (independent English benchmark questions).

\subsection{Hypotheses}
% Add directional hypotheses only if agreed before confirmatory analysis.
% Otherwise label analyses exploratory and avoid post-hoc hypotheses.

\section{Models and Datasets}

\subsection{Model and Tokenizer Selection}
% For every model record exact HF repository ID, immutable revision/commit,
% tokenizer class, vocabulary size, model size, context length, and access date.
% Reconcile the model panels in the existing branches before choosing the final panel.

\subsection{Dataset}

For our study, we use the BnMMLU: Measuring Massive Multitask Language Understanding in Bengali \cite{joy-shatabda-2026-bnmmlu} dataset. The dataset contains multiple-choice questions across the 23 domains, including science, humanities, mathematics, and general knowledge. It has primarily two categories: full, and hard.

The full category contains the samples which is simpler for large language models (LLMs) to answer, where as the hard category contains the samples which is difficult for the large language model to answer. Also, the hard category samples contains one extra feature named failure\_ratio, which highlights the failure on large language model during the study.

Since, our study address the cross-lingual study, that is why we translated the dataset into English. Below, we provide the full character of the dataset and translation procedures.

\subsubsection{Dataset Characteristics}

% Please add the following required packages to your document preamble:
\begin{table}[]
	\centering
	\begin{tabular}{@{}lllllll@{}}
		\cmidrule(r){1-5}
		Unique\_Serial & subject\_name            & question                                              & correct\_answer & options & question\_char\_count & faiture\_ratio \\ \cmidrule(r){1-5}
		1              & Algebra \& Number Theory & (x+a)/(x+1)(x-3) = B/(x+1) + 1/(x-3) হলে, A ও B এর মান- &                 &         &                       &                \\
		               &                          &                                                       &                 &         &                       &                \\
		               &                          &                                                       &                 &         &                       &                \\ \cmidrule(r){1-5}
	\end{tabular}
	\caption{}
	\label{tab:my-table}
\end{table}

\subsubsection{Translation Procedures of the Dataset}

\section{Tokenization Fragmentation Measurement}

\subsection{Text Unit Being Tokenized}
% State exactly whether the tokenizer receives question only, question plus options,
% or another fixed string. Use the same rule across models and languages.

\subsection{Tokenizer Execution and Reproducibility}
% Tokenizer-only measurement does not require model inference.
% Record tokenizer repository/revision, library version, special-token settings,
% add_special_tokens setting, and exact input construction.

\subsection{Word Counting and Normalization}
% Define word segmentation for Bangla and English, punctuation treatment,
% whitespace, digits, symbols, Unicode normalization, zero-word strings,
% and whether special tokens are excluded. These rules must be shared by all authors.

\subsection{Tokens per Word (TPW)}
% Primary proposed fragmentation measure. Confirm final counting rules before analysis.
\begin{equation}
	\mathrm{TPW}(x) =
	\frac{N_{\mathrm{tokens}}(x)}
	{N_{\mathrm{words}}(x)}.
	\label{eq:tpw}
\end{equation}
% Explain aggregation: item-level TPW first; then report distributions/summary
% statistics by model, language, and subject. Do not report only one global mean.

\subsection{Subword Token Retention Ratio (STRR)}
% Include only after Naeem/team verifies the exact definition from its source.
% State formula, numerator, denominator, token/word alignment method,
% treatment of whole-word tokens, and implementation tests.
% If the definition cannot be verified, omit STRR from primary analyses and
% explain this decision rather than inventing a formula.

\subsection{Worked Tokenization Examples}
% Present five manually checked Bangla-English pairs.
% For each: pair ID, exact input text, tokenizer/model revision, token pieces,
% token count, word count, TPW, and verified STRR (if retained).
% Do not include model answer/inference output in this tokenizer-only table.

% Suggested Table 2: five worked examples.
% Columns: Pair ID | Language | Text | Token pieces | #Tokens | #Words | TPW | STRR.

\section{Confidence and Calibration Methodology}
\subsection{Prompt and Multiple-Choice Scoring Protocol}
% Provide exact prompt template and chat template.
% Explain how A/B/C/D scores are extracted, normalized, and mapped to a prediction.
% Specify handling of multi-token answer labels, invalid outputs, ties, and missing scores.

\subsection{Confidence Definition}
% Define confidence as the probability assigned to the predicted answer, if that
% is the chosen protocol. State whether scores are logits, normalized probabilities,
% verbal self-reports, or another quantity. Do not mix these approaches.

\subsection{Accuracy}
\begin{equation}
	\mathrm{Accuracy} = \frac{1}{N}\sum_{i=1}^{N}
	\mathbb{1}(\hat{y}_i=y_i).
	\label{eq:accuracy}
\end{equation}
% Define the evaluation population and treatment of invalid/missing predictions.

\subsection{Brier Score}
\begin{equation}
	\mathrm{Brier} =
	\frac{1}{N}\sum_{i=1}^{N}\sum_{k=1}^{K}
	(p_{ik}-y_{ik})^2.
	\label{eq:brier}
\end{equation}
% State whether this multiclass convention (sum over classes) is used consistently.

\subsection{Expected Calibration Error (ECE)}
\begin{equation}
	\mathrm{ECE} =
	\sum_{b=1}^{B}\frac{|I_b|}{N}
	\left|\mathrm{acc}(I_b)-\mathrm{conf}(I_b)\right|.
	\label{eq:ece}
\end{equation}
% Freeze bin count, bin boundaries, endpoint convention, confidence definition,
% and empty-bin handling. Report bin count and sample size beside every ECE value.

\subsection{Reliability Diagrams}
% Explain binning and plot observed accuracy against mean confidence.
% Include the diagonal perfect-calibration reference line.
% Plot Bangla and English separately or clearly distinguish them.

\subsection{Calibration Variants and Data Leakage Prevention}
% If temperature scaling is used, specify the calibration/validation split,
% fitting procedure, and test split. Never fit temperature on the test set.
% Clearly separate raw (uncalibrated) and temperature-scaled results.

\section{Experimental Setup}
\subsection{Evaluation Pipeline}
% Describe code entry points, environment, package versions, hardware,
% precision, batch size, random seeds, caching, retries, and checkpoint/resume.
% Link the exact repository commit used for the final run.

\subsection{Primary and Secondary Outcomes}
% Name one primary fragmentation metric and one primary calibration outcome.
% Mark all other metrics and subgroup comparisons as secondary/exploratory.

\subsection{Statistical Analysis Plan}
% Crucial: the observational unit and dependence structure must be explicit.
% TPW is available per item and tokenizer; ECE is an aggregate metric, not an
% independent per-question label. Do not correlate item TPW with one global ECE
% value as if each question had its own ECE.
% Define the item-level calibration-error outcome or an appropriate grouped/
% hierarchical analysis; account for repeated items within models/subjects.
% Specify effect sizes, confidence intervals, clustering/mixed effects,
% multiple-comparison correction, missing-data handling, and sensitivity analyses.
% Avoid treating thousands of repeated questions as independent model-level replicates.

\subsection{Robustness and Sensitivity Checks}
% Examples if supported by data: question-only vs question-plus-options TPW;
% alternative word segmentation; paired-only subset; subject-level analysis;
% raw vs temperature-scaled confidence; alternative ECE bin counts.
% State which are planned, not all need to be run.

\section{Results}
\subsection{Dataset and Model Coverage}
% Report final sample counts, model revisions, paired coverage, and exclusions.

\subsection{Tokenization Fragmentation by Language and Model}
% Report TPW distributions (median, IQR, mean/SD if appropriate) by model/language.
% Add STRR only if definition and code have been verified.

\subsection{Accuracy and Calibration Performance}
% Report accuracy, Brier score, ECE, sample sizes, and uncertainty where available.
% Keep raw and temperature-scaled rows separate.

\subsection{Association Between Fragmentation and Calibration}
% Report the pre-specified association analysis, effect size, confidence interval,
% p-value where appropriate, and model/subject dependence handling.
% Distinguish confirmatory from exploratory findings.

\subsection{Robustness Checks}
% Report only analyses actually performed under the documented protocol.

\section{Discussion}
\subsection{Summary of Main Findings}
% Answer each RQ directly without overstating the evidence.

\subsection{Interpretation and Comparison with Prior Work}
% Explain how findings agree/disagree with reviewed work and why.

\subsection{Implications for Bangla LLM Evaluation}
% Discuss implications supported by results, not speculation.

\section{Limitations and Threats to Validity}
\subsection{Construct Validity}
% TPW/STRR may not capture every aspect of linguistic difficulty or tokenization quality.

\subsection{Internal Validity}
% Prompting, probability extraction, model versions, decoding, and translation choices.

\subsection{External Validity}
% Selected models, benchmark domains, Bangla varieties, and generalization to other tasks.

\subsection{Statistical Conclusion Validity}
% Sample dependence, number of models/clusters, uncertainty, multiple testing,
% and observational (non-causal) design.

\section{Conclusion and Future Work}
% Concisely answer the research question, state the contribution and limits,
% and identify the next experiment. Do not introduce new results here.

\begin{credits}
	\subsubsection{\ackname}
	% Add funding and acknowledgements only after author confirmation.

	\subsubsection{\discintname}
	% TODO: All authors must confirm the competing-interest declaration.
	The authors declare that they have no competing interests relevant to this work.
\end{credits}

% TODO: Add a data/code availability statement if required by the target venue.

\bibliographystyle{splncs04}
\bibliography{references}

\end{document}

% ---------------------------------------------------------------------------
% FIGURE AND TABLE PLAN (implementation checklist; add actual floats above)
%
% Figure 1. Study workflow: dataset -> paired language conditions -> tokenizer
% fragmentation -> model scoring/confidence -> calibration -> association analysis.
%
% Figure 2. Tokenization examples: visual comparison of token pieces for selected
% Bangla-English matched items (optional; may be replaced by Table 2).
%
% Figure 3. TPW distributions: box/violin plots by model and language.
%
% Figure 4. Reliability diagrams: confidence vs empirical accuracy for each
% language/model (avoid overcrowding; use panels if needed).
%
% Figure 5. Fragmentation-calibration association: scatter/partial-effect plot
% using the exact statistical unit and model described in the analysis plan.
%
% Table 1. Dataset composition, pairing, filtering, and final sample counts.
% Table 2. Five worked tokenizer examples with token pieces and metric calculations.
% Table 3. Model and tokenizer identities with immutable revisions.
% Table 4. TPW/STRR descriptive statistics by model and language.
% Table 5. Accuracy, Brier score, and ECE with N and bin count.
% Table 6. Primary association model: coefficient/effect, CI, p-value, controls.
% Table 7. Robustness/sensitivity analysis (only if conducted).
%
% Figures: export as PDF for vector plots where possible; use legible fonts,
% axis labels, units, captions, and consistent language/model ordering.
% Every figure/table must be cited in the text and have a self-contained caption.
% Never create a plot from placeholder or unverified values.
